\documentclass[10pt,a4paper]{amsart}
\usepackage[margin=19mm]{geometry}
\usepackage{amsmath,amssymb,mathtools}
\usepackage[scr=rsfs,cal=euler]{mathalfa}
\usepackage{xfrac}
\usepackage[unicode,hidelinks,bookmarks=false]{hyperref}
\newcommand{\ie}{i.e.\ }
\newcommand{\cf}{cf.\ }
\newcommand{\resp}{respectively\ }
\newcommand{\Subsection}{\subsection}
\providecommand{\nobreakdash}{\nobreak\hbox{-}}
\setlength{\emergencystretch}{2em}
\multlinegap0pt
\usepackage{bm}
\usepackage{bbm}
\usepackage{dsfont}
\usepackage{xspace}
\usepackage{cancel}
\usepackage{mathtools}
\usepackage{amsthm}
\makeatletter
\@ifundefined{lemma}{\newtheorem{lemma}{Lemma}}{}
\@ifundefined{theorem}{\newtheorem{theorem}{Theorem}}{}
\makeatother
\title[Age of Information Optimization for Status Updates in Integr]{Age of Information Optimization for Status Updates in Integrated Sensing and Communication Systems}
\author{Marco Zanni, Mohamad Assaad, Touraj Soleymani}
\date{Source: arXiv:2605.24714v1 (CC BY 4.0)}
\begin{document}
\maketitle

\noindent\emph{Source and licence.} Age of Information Optimization for Status Updates in Integrated Sensing and Communication Systems. Version arXiv:2605.24714v1 (CC BY 4.0), distributed under the Creative Commons Attribution 4.0 International licence, as recorded in the arXiv licence metadata for this exact version and shown on the abstract page https://arxiv.org/abs/2605.24714v1. Full rights notice: licenses/arxiv-2605.24714-CC-BY-4.0.txt.\par\smallskip

\noindent\textit{Editorial scope.} Counted unit: the Theorem whose source label is \texttt{th:policy-structure} in the article (source0.tex lines 281--292, source label \texttt{th:policy-structure}), delivered together with the article's own complete original proof of it (source0.tex lines 548--649, one \texttt{proof} environment of 102 lines). No mathematics is written, repaired or re-derived: statement and proof are the article's own text line for line. The proof is detached from the statement in the source: the article states the result at lines 281--292 and prints its proof at lines 548--649, and the proof environment's own header names the statement, so the association is the article's own. Required context retained uncounted: eq:problem-formulation (lines 233--236); lemma:ordered-cost-diagonal-zero (lines 446--448); lemma:delta-non-decreasing-in-a\_s (lines 522--524); lemma:delta-non-increasing-in-a\_b (lines 533--535). Every definition, display and label of those lines is retained unchanged. Omitted from this package and not redistributed: 1--232, 237--280, 293--445, 449--521, 525--532, 536--547, 650--1705. Nothing inside a retained excerpt is omitted or altered: every retained body is delivered exactly as the article's lines are written, so no omission receipt of any kind accompanies this package. Citations: the retained excerpts carry no citation command at all, so no bibliography entry of the article is transcribed and no reference is redistributed. Reference adaptation: none was needed and none was made; every reference inside the delivered unit resolves inside it, so no unresolved reference or citation is produced. Printed slips kept literal: no printed word, symbol, hypothesis or number of the counted statement or of its proof was altered. Layout and class: the article's own class is \texttt{IEEEtran} with packages including babel, graphicx, xcolor, amsmath, bm, amsfonts, amssymb, algorithm, algpseudocode, dsfont, amsthm, comment, cite, booktabs; the wrapper is the lane's pinned \texttt{amsart} editorial wrapper at 10pt on A4, which restates the theorem-like environments the retained text needs and adds the article's own macro definitions that the retained text needs (none), with extra wrapper packages (bm, bbm, dsfont, xspace, cancel, mathtools). The delivered numbering is the unit's own. Assets: no retained span contains a figure environment, a graphics-inclusion command, a PDF-page inclusion, a table or a TikZ picture; no image file is copied into this package, the package declares no asset dependency and no image file is redistributed with it. Licence: version arXiv:2605.24714v1 of this article is distributed under the Creative Commons Attribution 4.0 International licence, as recorded in the arXiv licence metadata for this exact version and shown on the abstract page https://arxiv.org/abs/2605.24714v1. A full rights notice is delivered at licenses/arxiv-2605.24714-CC-BY-4.0.txt. Characters: every retained excerpt is pure ASCII, so the delivered engine, XeLaTeX with its default Latin Modern text font, reports no missing character.
\begin{equation}
	\min_{\pi \in \mathcal P} \mathbb{E}\left[\sum_{k=0}^{\infty} \gamma^k g(S_k,u_k)\right],
	\label{eq:problem-formulation}
\end{equation}

\begin{lemma}\label{lemma:ordered-cost-diagonal-zero}
For every $\alpha\geq 1$, action $0$ is optimal on the diagonal state $(\alpha,\alpha)$.
\end{lemma}

\begin{lemma}\label{lemma:delta-non-decreasing-in-a_s}
For each fixed $\alpha^b$, the functions $\Delta^*_{01}(\alpha^m,\alpha^b)$, $\Delta^*_{02}(\alpha^m,\alpha^b)$, and $\Delta^*_{21}(\alpha^m,\alpha^b)$ are nondecreasing in $\alpha^m$.
\end{lemma}

\begin{lemma}\label{lemma:delta-non-increasing-in-a_b}
For each fixed $\alpha^m$, the functions $\Delta^*_{01}(\alpha^m,\alpha^b)$ and $\Delta^*_{02}(\alpha^m,\alpha^b)$ are nonincreasing in $\alpha^b$ for $1\leq \alpha^b\leq \alpha^m-1$.
\end{lemma}

\begin{theorem}\label{th:policy-structure}
The problem in \eqref{eq:problem-formulation} admits an optimal stationary deterministic policy $\pi^*$ with a switching-threshold structure. Specifically, there exist two functions $\tau_1, \tau_2 : \mathbb{N}\rightarrow\mathbb{N}_0\cup\{+\infty\}$ such that
\[
u^*(\alpha^m,\alpha^b)=
\begin{cases}
	\mathsf{sense,} & \text{if } \alpha^m\leq\tau_1(\alpha^b),\\
	\mathsf{joint,} & \text{if } \tau_1(\alpha^b)<\alpha^m\leq\tau_2(\alpha^b),\\
	\mathsf{comm,} & \text{if } \alpha^m>\tau_2(\alpha^b).\\
\end{cases}
\]
for all $(\alpha^m,\alpha^b)\in\mathcal{S}$, where $\tau_1(\alpha^b)$ is nondecreasing in $\alpha^b$ and $\tau_1(\alpha^b)\leq\tau_2(\alpha^b)$ for every $\alpha^b$.
\end{theorem}

\begin{proof}[Proof of Theorem \ref{th:policy-structure}]
For each state $(\alpha^m,\alpha^b)$, define
\[
u^*(\alpha^m,\alpha^b)=
\begin{cases}
0, & \begin{array}[t]{l}
\text{if } \Delta^*_{01}(\alpha^m,\alpha^b)\le 0\text{ and}\\
\Delta^*_{02}(\alpha^m,\alpha^b)\le 0,
\end{array}\\[0.4em]
1, & \begin{array}[t]{l}
\text{if } \Delta^*_{01}(\alpha^m,\alpha^b)>0\text{ and}\\
\Delta^*_{21}(\alpha^m,\alpha^b)\ge 0,
\end{array}\\[0.4em]
2, & \text{otherwise.}
\end{cases}
\]
This rule is optimal. Indeed, if $\Delta^*_{01}\le 0$ and $\Delta^*_{02}\le 0$, then $Q_0^*\le Q_1^*$ and $Q_0^*\le Q_2^*$, so action $0$ is optimal. If $\Delta^*_{01}>0$ and $\Delta^*_{21}\ge 0$, then $Q_1^*<Q_0^*$ and $Q_1^*\le Q_2^*$, so action $1$ is optimal. In all remaining cases action $2$ is optimal: if $\Delta^*_{01}\le 0$ and the first case fails, then necessarily $\Delta^*_{02}>0$, hence $Q_2^*<Q_0^*\le Q_1^*$; if $\Delta^*_{01}>0$ and the second case fails, then necessarily $\Delta^*_{21}<0$, hence $Q_2^*<Q_1^*<Q_0^*$.

Now fix $\alpha^b\in\mathbb N$ and define
\begin{align*}
\mathcal A_0(\alpha^b)
&:=\{\alpha^m\ge \alpha^b:\ \Delta^*_{01}(\alpha^m,\alpha^b)\le 0,\\
&\qquad \Delta^*_{02}(\alpha^m,\alpha^b)\le 0\},\\
\mathcal A_1(\alpha^b)
&:=\{\alpha^m\ge \alpha^b:\ \Delta^*_{01}(\alpha^m,\alpha^b)>0,\\
&\qquad \Delta^*_{21}(\alpha^m,\alpha^b)\ge 0\}.
\end{align*}
By construction, action $0$ is optimal on $\mathcal A_0(\alpha^b)$, action $1$ is optimal on
$\mathcal A_1(\alpha^b)$, and action $2$ is optimal on the complement of
$\mathcal A_0(\alpha^b)\cup \mathcal A_1(\alpha^b)$ in $\{\alpha^m\ge \alpha^b\}$.

By Lemma \ref{lemma:ordered-cost-diagonal-zero}, $\alpha^b\in\mathcal A_0(\alpha^b)$, so $\mathcal A_0(\alpha^b)$ is nonempty. Moreover, by Lemma \ref{lemma:delta-non-decreasing-in-a_s}, both $\Delta^*_{01}$ and $\Delta^*_{02}$ are nondecreasing in $\alpha^m$. Hence, if $\alpha^m\in\mathcal A_0(\alpha^b)$ and $\tilde\alpha^m$ satisfies $\alpha^b\le \tilde\alpha^m\le \alpha^m$, then
\begin{align*}
\Delta^*_{01}(\tilde\alpha^m,\alpha^b)
&\le \Delta^*_{01}(\alpha^m,\alpha^b)\le 0,\\
\Delta^*_{02}(\tilde\alpha^m,\alpha^b)
&\le \Delta^*_{02}(\alpha^m,\alpha^b)\le 0,
\end{align*}
so $\tilde\alpha^m\in\mathcal A_0(\alpha^b)$. Therefore $\mathcal A_0(\alpha^b)$ is an initial
segment of $\{\alpha^m\ge \alpha^b\}$.

Similarly, by Lemma \ref{lemma:delta-non-decreasing-in-a_s}, both
$\Delta^*_{01}$ and $\Delta^*_{21}$ are nondecreasing in
$\alpha^m$. Hence, if $\alpha^m\in\mathcal A_1(\alpha^b)$ and $\tilde\alpha^m\ge \alpha^m$, then
\begin{align*}
\Delta^*_{01}(\tilde\alpha^m,\alpha^b)
&\ge \Delta^*_{01}(\alpha^m,\alpha^b)>0,\\
\Delta^*_{21}(\tilde\alpha^m,\alpha^b)
&\ge \Delta^*_{21}(\alpha^m,\alpha^b)\ge 0,
\end{align*}
so $\tilde\alpha^m\in\mathcal A_1(\alpha^b)$. Therefore $\mathcal A_1(\alpha^b)$ is a terminal
segment of $\{\alpha^m\ge \alpha^b\}$.

We may thus define
\[
\tau_1(\alpha^b):=\sup \mathcal A_0(\alpha^b)\in \mathbb N_0\cup\{+\infty\},
\]
\[
\tau_2(\alpha^b):=
\begin{cases}
\inf \mathcal A_1(\alpha^b)-1, & \text{if } \mathcal A_1(\alpha^b)\neq\emptyset,\\
+\infty, & \text{if } \mathcal A_1(\alpha^b)=\emptyset.
\end{cases}
\]
Since $\mathcal A_0(\alpha^b)$ is an initial segment and $\mathcal A_1(\alpha^b)$ is a terminal
segment, we have
\[
\mathcal A_0(\alpha^b)=\{\alpha^m\ge \alpha^b:\alpha^m\le \tau_1(\alpha^b)\}
\]
\[
\mathcal A_1(\alpha^b)=\{\alpha^m\ge \alpha^b:\alpha^m>\tau_2(\alpha^b)\}.
\]
Because the two sets are disjoint, $\tau_1(\alpha^b)\le \tau_2(\alpha^b)$. Hence
\[
u^*(\alpha^m,\alpha^b)=
\begin{cases}
0, & \text{if } \alpha^m\le \tau_1(\alpha^b),\\
2, & \text{if } \tau_1(\alpha^b)<\alpha^m\le \tau_2(\alpha^b),\\
1, & \text{if } \alpha^m>\tau_2(\alpha^b).
\end{cases}
\]

It remains to prove that $\tau_1$ is nondecreasing in $\alpha^b$. Let
$\alpha^m\in\mathcal A_0(\alpha^b)$ with $\alpha^m\ge \alpha^b+1$. By
Lemma \ref{lemma:delta-non-increasing-in-a_b},
\begin{align*}
\Delta^*_{01}(\alpha^m,\alpha^b+1)
&\le \Delta^*_{01}(\alpha^m,\alpha^b)\le 0,\\
\Delta^*_{02}(\alpha^m,\alpha^b+1)
&\le \Delta^*_{02}(\alpha^m,\alpha^b)\le 0,
\end{align*}
so $\alpha^m\in\mathcal A_0(\alpha^b+1)$. Therefore
\[
\mathcal A_0(\alpha^b)\setminus\{\alpha^b\}\subseteq \mathcal A_0(\alpha^b+1).
\]
Since Lemma \ref{lemma:ordered-cost-diagonal-zero} gives
$\alpha^b+1\in\mathcal A_0(\alpha^b+1)$, and both sets are initial segments, it follows that
\[
\tau_1(\alpha^b)\le \tau_1(\alpha^b+1).
\]
Thus $\tau_1$ is nondecreasing in $\alpha^b$.
\end{proof}

\end{document}
